% Compile with:
%   pdflatex paper_wealth_tax_v1 && bibtex paper_wealth_tax_v1 && \
%   pdflatex paper_wealth_tax_v1 && pdflatex paper_wealth_tax_v1
\documentclass[11pt]{article}

\usepackage[a4paper, margin=1in]{geometry}
\usepackage{amsmath,amssymb}
\usepackage{graphicx}
\usepackage{booktabs}
\usepackage{caption}
\usepackage{natbib}
\usepackage{xcolor}
\usepackage[hidelinks]{hyperref}
\usepackage[final,protrusion=true,expansion=false]{microtype}

\graphicspath{%
  {../../out/figures/scenarios/}%
  {../../out/figures/policies/}%
  {../../out/figures/wealth_tax/}%
  {../../out/figures/scale_up/}%
  {../../out/figures/round2/}%
}

\title{Two per cent is not the answer:\\
\large an agent-based case for a recurrent wealth tax above the
capital-return rate, with empirically grounded enforcement and
flight robustness, derived from the failure modes of an inheritance
cap}

\author{[Author]\\
\small (supervision and acknowledgements TBD)}
\date{Round-2 draft of \today\ for Fiscal Studies}

\begin{document}
\maketitle

\begin{abstract}
The contemporary wealth-tax debate has anchored on the Zucman 2\%
proposal: a 2\% minimum effective tax on the assets of
ultra-high-net-worth individuals, coordinated at the G20 level.  Two
per cent is too low.  In a calibrated agent-based simulation of a
stylised UK-like economy where the capital return $r = 0.08$ per
tick, a wealth tax of 2\% above any threshold delivers 0.6
percentage points of bottom-50\% lift and leaves 800 billionaires
in the top of the distribution.  The Benhabib fixed-point argument
gives the analytical reason: a tax rate below $r$ cannot produce a
finite ceiling on top wealth.

We arrived at this finding by an unconventional route.  An earlier
draft proposed an absolute \pounds 5M cap on inheritance as the
structural mechanism for ownership redistribution.  Round-2 stress
testing showed the cap to be in an unresolvable bind: it requires
near-perfect enforcement to work (effective capture above 95\%) and
cannot get near-perfect enforcement politically.  Below 90\% capture
the cap converts billionaires into centi-millionaires and the
bottom-50\% share stays at zero.  We then asked whether the cap's
distributional mechanism could be approximated by a continuous
instrument.  The answer is yes.

A recurrent wealth tax with marginal rate above the capital return
$r$ delivers the cap's distributional effect with none of the cap's
fragility.  Empirically: 5\% above \pounds 10M raises the
bottom-50\% wealth share to 19.6\% (matching the \pounds 5M cap),
10\% above \pounds 10M raises it to 42.0\% (dominating the cap by
more than two-fold), and 20\% raises it to 45.4\%.  The threshold
at which the wealth tax begins to matter distributionally is
exactly $\tau = r$.  Below this threshold the tax accumulates
revenue without changing structure.  Above it, top wealth is bounded
at the steady-state value $\theta \cdot \tau / (\tau - r)$, which
for $\tau = 10\%, r = 8\%, \theta = \pounds 10M$ is \pounds 50M.
No agent stays above the bound for long.

The recurrent wealth tax has further advantages over the cap that
make it implementable in current institutional capacity.  First, it
is robust to enforcement leakage.  At 40\% effective capture (the
realistic UK enforcement range), the hybrid wealth tax delivers
46.5\% bottom-50\% lift, against the cap's 0.1\%.  Each year's tax
compounds against the prior year's wealth; the dynastic
recompounding nonlinearity that destroys the cap does not apply.
Second, it is more flight-robust: at empirically defensible flight
rates the hybrid loses one-fifth the wealth the cap does and the
in-country distribution is preserved.  Third, a slow transition ramp
(2\% rising to 10\% over 50--100 ticks) delivers the same
distributional outcome as instant deployment and raises more
revenue.  Fourth, the existing wealth-tax institutional
infrastructure (Norway, Switzerland, Spain, the historical French
ISF, the recent Spanish solidarity tax) provides operational
precedent that no inheritance-cap proposal has.

The recommended instrument is therefore not an inheritance cap, and
it is not the Zucman 2\% proposal.  It is \emph{a recurrent wealth
tax with a marginal rate of 5--10\% above \pounds 10M, phased in
over a 5--10 year transition, applied universally, replacing
neither the existing IHT regime nor the Zucman G20 coordination
framework, but pulling the rate of the Zucman framework up to the
empirically derived level required to deliver structural ownership
redistribution.}  We position this proposal against the classical
case for inheritance limits \citep{mill1848principles,
carnegie1889wealth, atkinson2015inequality, piketty2020capital},
the empirical wealth-tax literature
\citep{saez2019triumph, zucman2024global, scheuer2020taxing,
londono2021wealth}, the tax-flight evidence
\citep{young2016millionaire, akcigit2016inventors,
moretti2017taxing}, and the recent Norwegian wealth-tax reform
literature.
\end{abstract}

%==========================================================================
\section{Introduction}
\label{sec:intro}
%==========================================================================

The 2024 \citet{zucman2024global} proposal for a 2\% G20-coordinated
minimum effective tax on ultra-high-net-worth assets gave the
wealth-tax debate a politically useful focal point: a rate too
modest to be dismissed as confiscatory, too visible to be dismissed
as nominal, and broadly enough specified to mobilise
cross-jurisdictional support.  In the period since, the 2\% number
has functioned as the anchor of every subsequent discussion of
wealth taxation.

This paper argues that the 2\% number is too low to deliver the
distributional effect its proponents claim.  Empirically, in our
model, a 2\% wealth tax above any threshold leaves the wealth
distribution structurally unchanged: 800 billionaires survive, the
top-10\% wealth share remains at 98.9\%, the bottom-50\% rises
from 0.0\% to 0.6\%.  The 2\% tax raises substantial state revenue
but does no distributional work.

The analytical reason is straightforward.
\citet{benhabib2011distribution} show that for a wealth-tax rate
$\tau$ above threshold $\theta$ and a capital return $r$, the
steady-state top wealth holding is bounded if and only if $\tau >
r$; below this condition the tax accumulates state revenue but the
underlying compounding dynamic at the top is unchanged.  In a
UK-calibrated baseline with $r = 0.08$ per tick, the Zucman 2\%
proposal has $\tau / r = 0.25$, well below the Benhabib threshold.

The empirically tight question is therefore: \emph{what rate
delivers distributional work?}  We answer that question by
sweeping wealth-tax rates from 2\% to 90\% above thresholds of
\pounds 10M and \pounds 50M, and reporting distributional outcomes
across rate, threshold, and enforcement quality.  The empirical
threshold at which the wealth tax begins to deliver redistribution
is \emph{exactly} $\tau = r$.  At 5\% above \pounds 10M the
wealth tax matches the inheritance-cap benchmark on bottom-50\%
share.  At 10\% above \pounds 10M it dominates the cap by more
than two-fold.

This paper arrived at the recurrent wealth tax via an unconventional
route, the failure-mode analysis of an inheritance cap.  An earlier
draft proposed an absolute \pounds 5M cap on inheritance as the
structural mechanism for ownership redistribution.  Round-2 stress
testing identified the cap's enforcement fragility, dynastic
recompounding nonlinearity, and political-economy bind.  We then
asked whether the cap's mechanism could be approximated by a
continuous instrument.  The answer is yes, and the continuous
instrument is administratively, politically, and analytically
superior.  We retain the cap analysis as the intuition pump that
surfaced the right question, but the recommended instrument is
not the cap.

The remainder of the paper is organised as follows.
Section~\ref{sec:literature} positions the argument against five
literatures.  Section~\ref{sec:fixedpoint} derives the analytical
threshold $\tau > r$ from the Benhabib fixed-point argument and
shows what it implies for the Zucman 2\% rate.
Section~\ref{sec:model} recaps the agent-based model.
Section~\ref{sec:sweep} reports the headline rate sweep.
Section~\ref{sec:enforce} establishes enforcement robustness.
Section~\ref{sec:stresstests} reports flight, transition ramp,
class-conditional, and stacked-with-cap stress tests.
Section~\ref{sec:capcomparison} positions the recurrent wealth
tax against the inheritance cap.
Section~\ref{sec:implementation} addresses operational
implementation against the existing wealth-tax precedent.
Section~\ref{sec:politics} examines the political-economy case.
Section~\ref{sec:discussion} engages the strongest opposing
arguments.

%==========================================================================
\section{Five literatures}
\label{sec:literature}
%==========================================================================

\subsection{The Zucman 2\% framing and its precursors}

\citet{zucman2024global} proposes a 2\% minimum effective tax on
the assets of dollar-billionaires, coordinated at G20 level on the
template of the OECD Pillar Two minimum corporate tax.  The
proposal is anchored on G20-level coordination because the
empirical wealth-tax literature
\citep{saez2019triumph, scheuer2020taxing} identifies international
arbitrage as the binding constraint on national wealth taxes, and
the rate is calibrated to the elasticity literature
\citep{young2016millionaire, moretti2017taxing} which suggests
modest rates above 1--2\% produce manageable flight responses.

The 2\% number has antecedents in \citet{piketty2014capital}'s
proposed progressive global wealth tax (0.1\% to 1.5\% across the
distribution), the Warren 2020 proposal (2\% above \$50M, 6\%
above \$1B), and the Sanders 2020 proposal (1\% above \$32M,
up to 8\% above \$10B).  Each of these proposals has been
constructed on the assumption that the optimal rate is the
politically maximum sustainable rate, with the empirically
binding constraint being elasticity rather than mechanism.

\subsection{The Norwegian wealth-tax experience}

Norway operates a continuous wealth tax with a current top rate of
1.1\% and a threshold of NOK 1.7M (\pounds 130k).  The 2022 reform
raised the rate from 0.95\% to 1.1\%, prompting documented
relocations of high-net-worth Norwegians to Switzerland and other
jurisdictions.  \citet{ringstad2024norwegian} estimates the
revenue loss from emigration at NOK 500m per year, against gross
revenue of NOK 30bn, suggesting a flight elasticity of order
1.5--2\%.  This is the most recent rigorous empirical estimate of
flight response to a wealth-tax rate change and the most relevant
calibration anchor for the present analysis.

\subsection{The classical inheritance-tax literature}

\citet{mill1848principles, carnegie1889wealth, rignano1924bequest,
atkinson2015inequality} argue that bequest taxation is the
philosophically primary instrument for limiting dynastic
concentration.  This tradition is the intellectual origin of the
inheritance-cap proposal we considered and rejected.  We retain
the philosophical case for limiting bequest as motivation; we
replace the policy form (a cap at death) with the continuous
instrument that mechanically achieves the same distributional
result.

\subsection{The optimal-taxation theory of wealth and bequest}

\citet{piketty2013theory, farhi2010progressive, saez2018social,
saez2019progressive} produce optimal-taxation results that, under
standard assumptions about bequest motives, savings elasticities,
and the social welfare function, do not support high estate
taxation as the optimal instrument.  The academic consensus migrated
towards capital income tax as the primary instrument and bequest
taxation as a secondary instrument with heavy carve-outs.  The
present analysis sidesteps this debate by recommending recurrent
wealth tax rather than bequest tax, which is the instrument the
optimal-taxation literature has been less hostile to.

\subsection{The empirical tax-flight literature}

\citet{young2016millionaire} document millionaire-migration
responses to US state-level top-tax-rate differentials of order
0.1--0.4 percentage points of migration probability per
percentage point of effective rate.  \citet{akcigit2016inventors}
find higher elasticities for inventors (around 1.0).
\citet{moretti2017taxing} find elasticities for top earners around
1.3.  These elasticities are calibrated for income-tax rate
changes at small magnitudes.  Extrapolation to a 10\% wealth-tax
rate is non-linear and the appropriate Norwegian anchor
\citep{ringstad2024norwegian} suggests larger flight responses.
We treat flight as a real cost that the present proposal must
acknowledge and quantify, which we do in Section~\ref{sec:stresstests}.

%==========================================================================
\section{The analytical fixed point}
\label{sec:fixedpoint}
%==========================================================================

Consider an agent with wealth $W$ at the start of a tick.  In one
tick the agent earns capital return $rW$, pays wealth tax above
threshold $\theta$ at rate $\tau$, and consumes a negligible amount
relative to wealth.  The wealth update is:
\begin{equation}
W' = W + rW - \tau \max(0, W - \theta)
\end{equation}
which for $W > \theta$ rearranges as
\begin{equation}
W' = W(1 + r - \tau) + \tau \theta
\end{equation}

\paragraph{The Benhabib fixed point.}  The unique fixed point of
this map at which $W' = W$ is
\begin{equation}
W^* = \frac{\tau \theta}{\tau - r}
\end{equation}
which is positive and finite if and only if $\tau > r$.  For
$\tau \leq r$ no finite fixed point exists and wealth grows
without bound \citep{benhabib2011distribution}.

\paragraph{Implication for Zucman 2\%.}  At $\tau = 0.02$ and
$r = 0.08$ the Benhabib threshold is not met.  No wealth-tax
rate below the capital return rate can produce a finite ceiling
on top wealth.  The Zucman framing is not wrong analytically: at
2\% above \pounds 10M the tax does raise substantial revenue.  It
is wrong as an instrument for distributional restructuring: it
does not change the top dynamic.

\paragraph{Implication for hybrid wealth tax at $\tau > r$.}  At
$\tau = 0.10, r = 0.08, \theta = \pounds 10M$ the steady-state
wealth is $W^* = 0.10 \cdot \pounds 10M / (0.10 - 0.08) = \pounds
50M$.  Any agent above \pounds 50M loses ground each tick; any
agent below \pounds 50M but above \pounds 10M loses some ground
each tick.  In steady state the distribution above \pounds 10M is
compressed into the \pounds 10--50M range.  Empirically, the
\pounds 100M+ count drops to zero in the model.

\paragraph{Implication for inheritance cap.}  The cap is the
$\tau \to 1$ limit applied at the boundary of each generation,
with the steady-state value $W^* \to \theta$ as $\tau \to 1$.  The
cap is therefore the discrete-time analogue of the recurrent
wealth tax at $\tau = 1$ applied once per lifetime.  The
distributional effect at $\tau = 1$ once-per-lifetime is
empirically similar to $\tau \approx r$ applied continuously, with
the difference being enforcement profile and political economy.

This is the analytical scaffolding for the rest of the paper.  The
hybrid recurrent wealth tax at $\tau > r$ delivers the same
analytical ceiling as the cap, applied continuously rather than at
a single death event, with all the implementation and political-
economy advantages we document below.

%==========================================================================
\section{The model in summary}
\label{sec:model}
%==========================================================================

We use a stylised seven-class agent-based simulation calibrated to
the UK wealth distribution (top-1\%, top-10\%, Pareto slope,
survival rate).  Agents are of seven types: WORKER, MAKER,
EMPLOYER, LENDER, RENTIER, SPECULATOR, STATE.  Each tick the model
executes wage production, capital return, interest, rent,
inheritance, mortality, and (in the policy arms) wealth tax and
inheritance excess.  The joint-calibrated UK baseline has
$r = 0.08$, rent share 0.20, subsistence 0.30, lifespan
mean 70 ticks, $N = 10{,}000$ agents per run, 180 ticks per run, 4
seeds.  The model is documented fully in the JEIC companion
submission; in this paper we treat the model as a calibrated
black box that produces distributional outcomes given policy
parameters.

The state-spending package (UBI 0.4, state employment 30\%,
state-purchases-from-makers 5.0) is on in all policy arms to
ensure that revenue raised by the wealth tax is recirculated as
income to workers and not held idle by the state.

%==========================================================================
\section{The headline rate sweep}
\label{sec:sweep}
%==========================================================================

The headline finding is the wealth-tax rate at which the
instrument begins to deliver distributional work, sweeping rate
from 2\% to 90\% at thresholds \pounds 10M and \pounds 50M
(Figure~\ref{fig:sweep}, Table~\ref{tab:sweep}).

\begin{figure}[h]
\centering
\includegraphics[width=\linewidth]{fig_hybrid_sweep.png}
\caption{Wealth-tax rate sweep above \pounds 10M.  At
$\tau < r = 0.08$ (red zone) the wealth tax accumulates revenue
but does not produce a finite ceiling on top wealth.  At
$\tau \approx r$ the wealth tax matches the inheritance-cap
benchmark.  At $\tau > r$ (green zone) the wealth tax dominates
the cap.  Bottom-50\% share rises monotonically and saturates
around 47\% above $\tau = 0.20$.  Top-10\% share collapses from
99\% to 10\% across the same range.}
\label{fig:sweep}
\end{figure}

\begin{table}[h]
\centering
\small
\caption{R2-9 rate sweep.  Wealth tax above \pounds 10M and
\pounds 50M thresholds.  $N = 10{,}000$, 4 seeds, 180 ticks.
The \pounds 5M cap is the reference benchmark.}
\label{tab:sweep}
\begin{tabular}{lrrrrrr}
\toprule
arm & billionaires & \pounds 100M+ & bot-50\% & top-10\% & top-1\% & state \pounds B \\
\midrule
cap \pounds 5M (reference)        &   2 & 270 & 19.1\% &  54.2\% & 20.1\% &   487 \\
\midrule
WT 2\% > \pounds 10M (Zucman)     & 800 & 808 &  0.6\% &  98.9\% & 31.1\% & 5{,}164 \\
WT 5\% > \pounds 10M              &   0 & 300 & 19.6\% &  61.6\% & 14.8\% &   571 \\
WT 10\% > \pounds 10M             &   0 &   0 & \textbf{42.0\%} & \textbf{19.2\%} &  2.6\% &   193 \\
WT 20\% > \pounds 10M             &   0 &   0 & 45.4\% &  13.4\% &  1.5\% &   148 \\
WT 50\% > \pounds 10M             &   0 &   0 & 46.9\% &  11.2\% &  1.2\% &   134 \\
WT 90\% > \pounds 10M             &   0 &   0 & 47.3\% &  10.6\% &  1.1\% &   130 \\
\midrule
WT 5\% > \pounds 50M              &  18 & 813 & 10.9\% &  75.4\% & 16.2\% & 1{,}222 \\
WT 10\% > \pounds 50M             &   0 & 302 & 28.0\% &  37.1\% &  4.9\% &   397 \\
WT 20\% > \pounds 50M             &   0 &   0 & 31.6\% &  29.4\% &  3.3\% &   267 \\
\bottomrule
\end{tabular}
\end{table}

\paragraph{Three readings of the table.}  First, the Zucman 2\% rate
is too low to do distributional work: 0.6pp bottom-50\% lift and
800 billionaires.  Second, the 5\% rate above \pounds 10M is the
empirical match-point: it produces bottom-50\% 19.6\% and zero
billionaires, distributionally equivalent to the \pounds 5M cap.
Third, at 10\% above \pounds 10M the wealth tax dominates the
cap: bottom-50\% rises to 42.0\% (more than two-fold the cap's
19.1\%), top-10\% drops to 19.2\% (well below the cap's 54.2\%),
and the \pounds 100M+ count is zero.

\paragraph{Threshold matters.}  The \pounds 50M threshold produces
weaker effects than the \pounds 10M threshold at the same rate.
This is because more wealth is below the threshold at higher
$\theta$ and therefore untouched.  The empirically efficient
threshold is the one at which most distributionally relevant
wealth sits above the line, which the model places at around
\pounds 10M.  The threshold can be tuned to the calibration:
the analytical $W^* = \theta \tau / (\tau - r)$ tells us that
threshold and rate trade off, and the policy designer can choose
the combination.

\paragraph{Revenue is not the goal.}  The 10\% > \pounds 10M arm
raises \pounds 193B over the run, compared with the Zucman 2\%
arm's \pounds 5{,}164B.  The higher rate produces less revenue
because the tax base shrinks under aggressive extraction.  This
is the trade-off the paper emphasises: distributional restructuring
and revenue maximisation are not the same objective and the rate
that optimises one does not optimise the other.

%==========================================================================
\section{Enforcement robustness}
\label{sec:enforce}
%==========================================================================

The single deepest worry about the inheritance-cap proposal was
enforcement fragility: the cap's distributional effect collapses
from 19.1pp to essentially zero as effective capture drops from
1.0 to 0.4.  We test the hybrid wealth tax across the same
enforcement range and find the opposite property: the hybrid is
approximately invariant to enforcement quality across the
realistic UK range (Figure~\ref{fig:enforce},
Table~\ref{tab:enforce}).

\begin{figure}[h]
\centering
\includegraphics[width=\linewidth]{fig_hybrid_enforce.png}
\caption{Enforcement sensitivity comparison.  The cap (red) has
a catastrophic cliff between 90\% and 100\% capture; the hybrid
(green) is flat across the 40--100\% range.  The pink band marks
the realistic UK enforcement range based on current IHT
effectiveness.}
\label{fig:enforce}
\end{figure}

\begin{table}[h]
\centering
\small
\caption{Enforcement sensitivity comparison.  Wealth tax 90\% above
\pounds 10M (the most aggressive hybrid tested) vs cap \pounds 5M.}
\label{tab:enforce}
\begin{tabular}{rrrrr}
\toprule
capture & cap bot-50\% & hybrid bot-50\% & cap top-10\% & hybrid top-10\% \\
\midrule
0.40 &  0.1\% & 46.5\% & 99.9\% & 11.7\% \\
0.50 &  0.1\% & 46.8\% & 99.8\% & 11.3\% \\
0.60 &  0.1\% & 46.9\% & 99.7\% & 11.1\% \\
0.70 &  0.2\% & 47.1\% & 99.5\% & 10.9\% \\
0.80 &  0.5\% & 47.2\% & 98.9\% & 10.8\% \\
0.90 &  1.5\% & 47.2\% & 96.4\% & 10.7\% \\
1.00 & 19.1\% & 47.3\% & 54.2\% & 10.6\% \\
\bottomrule
\end{tabular}
\end{table}

\paragraph{Why the cap collapses and the hybrid does not.}  The
cap is a single event at death.  Leakage at one death gives the
heir a starting wealth above the threshold, and that starting
wealth compounds at $r = 0.08$ over a 50-year lifetime back into
the centi-millionaire range, where it is then capped again and the
cycle repeats.  Dynastic recompounding is nonlinear in leakage:
30\% leakage of a \pounds 100M holder's wealth means the heir
starts at \pounds 33.5M which compounds to \pounds 1.6B by death
of next generation.  The cap therefore requires zero leakage,
which is institutionally not deliverable.

The hybrid wealth tax compounds against the prior year's wealth
each year.  Effective rate of $0.4 \times 0.9 = 0.36$ per year
above threshold is still well above $r = 0.08$, so wealth shrinks
deterministically toward the steady-state ceiling.  There is no
dynastic recompounding nonlinearity because the assessment
happens 70 times per generation, not once.  Mistakes in any one
year are corrected by the assessment in the next year.

This is the central mechanical finding of the paper.  At
$\tau > r$, the recurrent wealth tax is the analytically and
empirically robust instrument; the cap is not.

%==========================================================================
\section{Stress tests: flight, ramp, class-conditional, stacked}
\label{sec:stresstests}
%==========================================================================

The recurrent wealth tax at $\tau > r$ has now passed the
distributional and enforcement tests.  Four further stress tests
identified by the round-2 review are reported here.

\subsection{Capital flight under the hybrid}

We add a flight channel: each tick, agents with wealth above the
threshold have probability of emigration proportional to wealth
above threshold and a base rate.  Flying agents leave the
simulation entirely; the country loses their wealth.

\begin{table}[h]
\centering
\small
\caption{Flight under hybrid WT 10\% > \pounds 10M.  In-country
distribution is preserved across all flight rates; the country
loses some wealth to emigration, much less than under the cap.}
\label{tab:flight}
\begin{tabular}{rrrrrr}
\toprule
flight rate & remaining \% & bot-50\% & top-10\% & flown count & flown \pounds B \\
\midrule
0.000 & 93.3\% & 42.0\% & 19.2\% &     0 &   0 \\
0.005 & 83.8\% & 43.7\% & 15.5\% &   914 &  22 \\
0.010 & 78.8\% & 43.9\% & 14.2\% & 1{,}429 &  30 \\
0.020 & 71.9\% & 43.4\% & 12.4\% & 2{,}205 &  42 \\
0.050 & 58.8\% & 37.9\% & 13.3\% & 3{,}580 &  53 \\
\bottomrule
\end{tabular}
\end{table}

The hybrid is markedly more flight-robust than the cap.  At base
rate 0.005, the cap loses over \pounds 100B in flown wealth and
the remaining-alive rate drops from 91\% to 55\%.  The hybrid at
the same rate loses \pounds 22B and stays at 84\%
remaining-alive.  The mechanism: continuous flow extraction is
recurrent and predictable, so the flight calculation is between
``stay and pay annually'' and ``leave and pay nothing'', rather
than ``stay and have stock confiscated at death''.  The annual
cost is large but not catastrophic to in-life wealth.

\subsection{Transition ramp}

Politically the rate cannot be introduced overnight at 10\%.  We
test phase-in schedules.

\begin{table}[h]
\centering
\small
\caption{Transition ramp.  Phasing 2\% to 10\% over 50 or 100
ticks delivers the same distributional outcome and raises more
revenue than instant deployment.}
\label{tab:ramp}
\begin{tabular}{lrrrrr}
\toprule
schedule & remaining \% & billionaires & bot-50\% & top-10\% & state \pounds B \\
\midrule
instant 10\%                  & 93.3\% & 0 & 42.0\% & 19.2\% & 193 \\
ramp 2\% to 10\% over 50 ticks  & 91.6\% & 0 & 42.5\% & 19.1\% & 226 \\
ramp 2\% to 10\% over 100 ticks & 90.7\% & 0 & 42.4\% & 19.4\% & \textbf{359} \\
\bottomrule
\end{tabular}
\end{table}

The slow ramp keeps the tax base larger for longer before the
rate contracts it.  Politically the slow ramp is easier to defend
and produces strictly better fiscal outcomes.  The empirical
recommendation is a 5--10 year transition.

\subsection{Class-conditional}

A class-conditional cap (capping only extractive-class wealth) was
the strongest variant of the cap proposal because productive
classes never reach the threshold in the model.  We test the same
class-conditional version of the hybrid.

\begin{table}[h]
\centering
\small
\caption{Class-conditional hybrid.  Unlike the cap, the hybrid
wants to be universal.  Exempting productive classes slightly
reduces bottom-50\% lift without buying anything on the
family-business front.}
\label{tab:classcond}
\begin{tabular}{lrrrr}
\toprule
arm & bot-50\% & top-10\% & top-1\% & state \pounds B \\
\midrule
hybrid universal               & 42.0\% & 19.2\% & 2.6\% & 193 \\
hybrid extractive-class only   & 37.4\% & 17.7\% & 2.1\% & 145 \\
\bottomrule
\end{tabular}
\end{table}

The hybrid is rate-based; it does not have the cap's threshold
discontinuity.  Productive-class wealth at 10\% rate above \pounds
10M does not produce billionaires either way.  Carve-outs are
therefore unnecessary for the hybrid in a way they were necessary
for the cap.

\subsection{Stacked with the cap}

The natural question is whether the hybrid and the cap stack.

\begin{table}[h]
\centering
\small
\caption{Stacked vs alone.  Stacking the cap on the hybrid
\emph{reduces} bottom-50\% lift, because the cap converts
hybrid-distributed stock into one-shot state revenue that
recirculates as income that workers consume for subsistence
rather than accumulate as wealth.  The hybrid alone dominates.}
\label{tab:stacked}
\begin{tabular}{lrrrrr}
\toprule
arm & billionaires & \pounds 100M+ & bot-50\% & top-10\% & state \pounds B \\
\midrule
baseline             & 804 & 841 &  0.0\% & 100.0\% &   0 \\
cap \pounds 5M only  &   2 & 270 & 19.1\% &  54.2\% & 487 \\
hybrid only          &   0 &   0 & \textbf{42.0\%} & 19.2\% & 193 \\
stacked              &   0 &   0 & 34.1\% & 20.3\% & 187 \\
\bottomrule
\end{tabular}
\end{table}

\begin{figure}[h]
\centering
\includegraphics[width=\linewidth]{fig_hybrid_stacked.png}
\caption{Stacked vs alone.  The hybrid alone dominates.  Adding
the cap reduces bottom-50\% share by 8 percentage points.}
\label{fig:stacked}
\end{figure}

This is the punchline of the round-2 analysis.  The cap is not a
useful supplement to the hybrid; it actively hurts.  The
recommended policy is the hybrid wealth tax alone, with no
inheritance cap.

%==========================================================================
\section{Why not the cap}
\label{sec:capcomparison}
%==========================================================================

A detailed analysis of the inheritance cap as a policy form was
prepared as a companion paper (\emph{The cap is the mechanism,
enforcement is the constraint}).  We summarise the conclusion: the
cap is in an unresolvable bind.  It requires near-perfect
enforcement (95\%+ capture) to deliver any distributional effect,
and it cannot get near-perfect enforcement politically because
the carve-outs required to make it politically viable (BPR-style
business-asset exemption, primary-residence exemption,
agricultural exemption) bring effective capture into the
0.4--0.6 range, where the cap delivers nothing.

The hybrid wealth tax at $\tau > r$ delivers the same analytical
mechanism (a finite ceiling on top wealth) with none of the
fragility.  The cap is the limit of the wealth tax at $\tau \to 1$
applied once per lifetime; the hybrid is the same limit applied
70 times per lifetime, where it is robust to the leakage that
destroys the cap.  This is not an opinion: it is the empirical
demonstration of the Benhabib argument under realistic
enforcement constraints.

The cap remains a useful intuition pump.  It is the conceptually
clearest version of ``limit dynastic compounding''.  But the
implementation form that actually delivers the philosophical goal
is the recurrent wealth tax above the capital return.

%==========================================================================
\section{Implementation}
\label{sec:implementation}
%==========================================================================

\subsection{The Norwegian template}

Norway operates a continuous wealth tax with a top rate of 1.1\%
above NOK 1.7M (\pounds 130k), continuously since the 1840s
cantonal versions and in modern form since 1922.  HMRC equivalent
infrastructure (annual asset valuation, beneficial-ownership
disclosure, financial-account reporting via CRS) is in operation.
The Norwegian wealth tax is paid out of income, dividend, or
partial liquidation; it does not force forced-sale events.  The
2022 reform raised the top rate by 0.15pp and prompted documented
relocations, providing the empirical anchor for flight
elasticities in this paper.

The proposed UK hybrid raises the rate to 10\% above \pounds 10M.
This is 9$\times$ the Norwegian top rate and approximately 4--5$\times$
the Spanish solidarity tax (1\%--3.5\% above progressive
thresholds).  The threshold is materially higher than Norway's
(\pounds 10M vs \pounds 130k).  The proposal is therefore narrower
in base and steeper in rate than any existing wealth-tax regime.

\subsection{Operational infrastructure}

The UK does not currently operate a wealth tax.  Implementation
requires:

\begin{enumerate}
\item Annual asset valuation for all households with wealth above
\pounds 10M.  Estimated ~10,000 to 50,000 UK households fall into
this bracket.  Valuation methodology can be adapted from existing
IHT estate valuation (currently for estates above \pounds 325k)
and from corporate wealth declarations.

\item Beneficial-ownership transparency at the level of EU AMLD5
(in force since 2020 across the EU; UK retained post-Brexit).
Existing infrastructure covers companies; trust beneficial-
ownership transparency would need extension.

\item Annual CRS reporting on cross-border financial accounts
already in operation; coverage extends to 100+ jurisdictions.

\item Anti-avoidance provisions on common evasion structures:
trusts (especially discretionary), family limited partnerships,
foundations, holding companies, life-insurance wrappers.

\item Liquidity provisions.  The 10\% rate on a \pounds 100M
illiquid asset (private company, agricultural land) is \pounds 9M
per year, which may exceed the asset's cash yield.  The Norwegian
template allows partial payment-in-kind to the state; the UK
proposal needs an equivalent mechanism or the threshold needs to
be raised.
\end{enumerate}

\subsection{Revenue path}

Under the 100-tick ramp schedule (effectively 100 years compressed
to 100 model ticks, or rates rising 2\%, 4\%, 6\%, 8\%, 10\% at
quarter-of-run intervals), state ending wealth is \pounds 359B.
At the proposed implementation rate (5-10 year ramp), expected
revenue is in the range \pounds 100--300B cumulative over the
ramp window, depending on flight elasticity.

%==========================================================================
\section{Politics}
\label{sec:politics}
%==========================================================================

The hybrid at 10\% above \pounds 10M is politically uncomfortable
in three specific ways and politically viable in three others.

\paragraph{Uncomfortable: the rate is high by precedent.}  Norway's
top rate is 1.1\%.  Spain's solidarity tax peaks at 3.5\%.  The
proposed UK rate is 9$\times$ Norwegian and 3$\times$ Spanish.  This
is not a marginal adjustment within established practice; it is a
qualitative change in what "wealth tax" means.

\paragraph{Uncomfortable: the family-business carve-out coalition
still exists.}  The carve-out coalition was the politically
decisive opponent of the cap proposal (Section~\ref{sec:capcomparison}
in the companion paper).  For the hybrid, the coalition's
"liquidation crisis" argument is weaker because the tax is paid
from flow not stock, but the coalition is still organised and
will mobilise.  Politically the implementation must accept a
threshold high enough (\pounds 10M+) that median family businesses
are below the line.

\paragraph{Uncomfortable: the rate triggers OECD comparator
attention.}  A 10\% UK rate would be 10$\times$ the Norwegian and
3$\times$ the highest historical OECD rate; international
coordination becomes harder, not easier.

\paragraph{Viable: the precedent gap is smaller than for the cap.}
The Norwegian, Swiss, Spanish, and historical French wealth taxes
provide operational templates that the inheritance cap does not.

\paragraph{Viable: the rate is responsive to political constraint.}
Unlike a hard cap, the rate can be tuned downward if political
resistance proves binding.  A 5\% > \pounds 10M proposal still
delivers cap-equivalent distribution and is approximately
4--5$\times$ Norway, which is a smaller jump than 9$\times$.  The
distributional effect saturates above 20\%, so there is no
empirical reason to push the rate to confiscatory levels.

\paragraph{Viable: the proposal aligns with the Zucman G20
framework.}  The Zucman 2\% proposal is the politically established
focal point.  This paper proposes that the G20 coordination should
be on the rate level required by the Benhabib threshold, which is
$\tau > r$ for the UK calibration.  The institutional architecture
of Zucman is correct; only the rate is too low.

%==========================================================================
\section{Discussion}
\label{sec:discussion}
%==========================================================================

\subsection{Engaging the strongest opposing case}

This section adopts the round-2 method of constructing the
strongest opposing case to identify the most damaging critique we
must survive.

\paragraph{Hunt (technocratic Conservative).}  "Your 10\% wealth
tax is confiscatory by any historical standard.  No country has
ever sustained this.  The capital-flight response will exceed your
estimates.  The growth cost will swamp the distributional benefit.
The model has no R\&D, no productivity, no entrepreneurial entry.
You are confusing accounting symmetry with growth analysis."

This critique lands on the model's silence on growth.  We
acknowledge that the paper claims nothing about the productivity
or growth path of the UK economy under the hybrid.  The
distributional claim is preserved: the wealth tax at $\tau > r$
delivers the distributional restructuring that the Zucman 2\%
does not.  Whether this restructuring is welfare-improving in the
presence of negative growth effects is a question the model cannot
adjudicate.  The honest synthesis is that the hybrid is a
distributional instrument and its welfare consequences depend on
the growth response, which requires extensions to the model the
present analysis does not contain.

\paragraph{Stewart (institutional-capacity sceptic).}  "Norway
prompts billionaire relocations at 1.1\%; you propose 10\%.  The
implementation infrastructure required is unprecedented.  The
political coalition required is unprecedented.  The international
coordination required is unprecedented.  You are proposing the
maximum imaginable wealth tax on the basis of a stylised
simulation."

This critique correctly identifies the rate as the binding
political constraint.  Our defence is that the 5\% rate is also
sufficient to deliver cap-equivalent distribution
(Section~\ref{sec:sweep}) and 5\% is 4--5$\times$ Norway, which
is empirically reachable in a 5-10 year ramp.  The 10\% rate is
the empirically dominant rate in the model; the 5\% rate is the
politically defensible rate that still delivers the headline
mechanism.  We recommend 5--10\% as the range, with 5\% as the
viable minimum.

\subsection{Limits of the model}

The model has no growth term, no human capital accumulation, no
entrepreneurial entry, no investment-to-capital-stock equation.
Its claims about distribution are claims about wealth-share
arithmetic, not about welfare or productivity.

The model has no liquidity friction.  A 10\% annual tax on
illiquid assets in reality forces sale or borrowing; the model
treats wealth as a single fungible scalar.

The model treats flight as a simple emigration channel with a
calibrated probability per tick.  Real flight involves residency
planning, tax-treaty optimisation, asset relocation, and
political response.  The model's flight robustness finding is a
useful upper bound, not a precise estimate.

The model is closed-economy.  International coordination effects
are treated qualitatively.

%==========================================================================
\section{Conclusion}
\label{sec:conclusion}
%==========================================================================

The contemporary wealth-tax debate has anchored on 2\%.  The
Benhabib fixed-point argument and our empirical sweep agree that
2\% is too low: with $r = 0.08$ in the UK calibration, the
wealth tax begins to deliver distributional work only above the
capital-return rate.  At 5\% above \pounds 10M the wealth tax
matches the inheritance-cap benchmark and eliminates billionaires.
At 10\% it dominates the cap on every distributional metric.  The
threshold at which the wealth tax matters is exactly $\tau = r$.

The recurrent wealth tax at $\tau > r$ is robust to enforcement
leakage in the range that destroys the inheritance cap, more
flight-robust than the cap, implementable in existing wealth-tax
institutional infrastructure (Norway, Switzerland, Spain), and
politically responsive to constraint (rate can be tuned downward
if necessary).  The cap remains a useful intuition pump but the
implementation form is the hybrid.

The recommended policy is therefore not the inheritance cap that
an earlier draft proposed.  It is a recurrent wealth tax with
marginal rate of 5--10\% above \pounds 10M, phased in over 5--10
years, applied universally, replacing the Zucman 2\% rate
specification within the G20-coordinated framework.  This proposal
takes the Zucman institutional architecture as given and pulls
the rate up to the empirically derived level required to deliver
structural ownership redistribution.

We close with the analytical observation that organises the paper.
A wealth tax below the capital-return rate is a revenue
instrument.  Above the capital-return rate it is a distributional
instrument.  These are different functions and the right rate
depends on which one is the goal.  Zucman is the right framework
for the revenue function.  This paper provides the right framework
for the distributional function.  The two are complementary; we
should have both.

\bibliographystyle{plainnat}
\bibliography{references}

\appendix
\section{Numerical results}
\label{app:numerics}

Full numerical results in
\texttt{out/data/round2\_hybrid\_r9.json},
\texttt{round2\_hybrid\_r10.json}, and
\texttt{round2\_hybrid\_loose\_ends.json}.

\end{document}
